\section{The resonant cell law}
\label{sec:resonant}

Throughout this section, $h>0$ and $E\in(-4-h,4+h)$ are fixed.
All choices may depend on these two parameters.  A polynomial bound in
$L$ means $C(1+L)^C$, with constants independent of the exterior system,
its dimension, and the artificial-probe mask.  The construction below
specifies an auxiliary law on one cell.  Its incorporation into the true
product law is deferred to \Cref{hier:sampling}.

\subsection{Clipping the corners}

Fix an integer margin $b\ge3$.  In a tile of odd side $b_l=3^l$, use
coordinates centered at its central site and put
\begin{equation}
 L=\frac{b_l-1}{2}-b,\qquad
 V=\{-L,\ldots,L\}^2\setminus\{(\pm L,\pm L)\},\qquad
 J=V\setminus\{0\}.
 \label{cell:sets}
\end{equation}
Only sufficiently large $l$ are used.  Let
\[
 B=\{x\notin V:\dist_1(x,V)=1\},
 \qquad p_x=\one_{\{|x|_1=1\}}\quad(x\in J).
\]
Thus $p$ is the adjacency column from the root $0$ to $J$.  Let
$M:\R^B\longrightarrow\R^J$ be the adjacency matrix from $B$ to $J$;
when used on $V$, it has an additional zero row at the root.  Define
\begin{equation}
 Z=\{x\in J:M_{x,\cdot}\ne0\},\qquad
 Z^+=\{x\in J:\dist_1(x,Z)\le1\},\qquad
 P=\operatorname{proj}_{(\ker M)^\perp}.
 \label{cell:boundary-spaces}
\end{equation}
We identify a projection with its range when specifying a matrix domain.
The cell and its corner coupling are shown in \cref{fig:clipped-cell}.

\begin{lemma}[The clipped boundary]\label{cell:geometry}
The set $Z$ consists of the four sides of the square with their corners
removed.  Every site of $Z$ has an inward nearest neighbor in $J\setminus
Z$.  The restriction $M:P\longrightarrow\R^Z$ is an isomorphism whose
singular values lie in $[1,\sqrt3]$.  In particular,
\[
 \norm{M}\le\sqrt3,\qquad
 \norm{Mw}\ge\norm{w}\quad(w\in P).
\]
There is an orthogonal decomposition $P=P_0+Q$: $P_0$ is the coordinate
projection onto the single side sites of $B$, and $Q$ is the direct sum
of four two-dimensional corner planes.
\end{lemma}

\begin{proof}
At the positive corner the relevant sites of $B$, in their fixed order,
are
\[
 e_1=(L+1,L-1),\qquad e_2=(L-1,L+1),\qquad e_3=(L,L).
\]
Here and below the same symbols also denote the corresponding coordinate
vectors.  The two rows indexed by $(L,L-1)$ and $(L-1,L)$ give the block
\begin{equation}
 \begin{pmatrix}1&0&1\\0&1&1\end{pmatrix}.
 \label{cell:corner-matrix}
\end{equation}
Its row space is $\operatorname{span}\{e_1+e_3,e_2+e_3\}$. We use the
orthonormal basis
\begin{equation}
 n=\frac{(-1,-1,1)}{\sqrt3},
 \qquad
 r_1=\frac{(1,-1,0)}{\sqrt2},\qquad
 r_2=\frac{(1,1,2)}{\sqrt6}.
 \label{cell:corner-basis}
\end{equation}
Here $n$ spans the kernel, and $r_1,r_2$ span the corner plane. The two
nonzero singular values are $1$ and $\sqrt3$. Use the
reflected geometric configuration and the same ordered basis coefficients
at the other three corners.  Every remaining nonzero row of $M$ contains
one entry, equal to $1$, in a distinct single side coordinate.  This
proves the asserted decomposition and singular-value bounds.  In fact
$|B|=8L$, $|Z|=8L-4$, and $\dim P_0=8L-12$.

For example, an inward neighbor of $(L,y)\in Z$ is $(L-1,y)$, which is
not in $Z$; the other sides follow by reflection.  These choices remain
valid next to a clipped corner.  All connections from $V$ to its
complement pass through $B$, and $V\cup B$ lies inside the tile because
of the fixed margin $b$.
\end{proof}

\begin{figure}[tb]
\centering
\begin{tikzpicture}[font=\small,>=Latex]
 \begin{scope}[x=0.48cm,y=0.48cm]
  \draw[dashed,gray] (-5.6,-5.6) rectangle (5.6,5.6);
  \fill[gray!12] (-4,-3)--(-3,-4)--(3,-4)--(4,-3)--
   (4,3)--(3,4)--(-3,4)--(-4,3)--cycle;
  \draw[thick] (-4,-3)--(-3,-4)--(3,-4)--(4,-3)--
   (4,3)--(3,4)--(-3,4)--(-4,3)--cycle;
  \foreach \j in {-3,-2,-1,0,1,2,3} {
   \foreach \sgn in {-1,1} {
    \fill (\sgn*4,\j) circle (0.07);
    \fill (\j,\sgn*4) circle (0.07);
    \draw[fill=white] (\sgn*5,\j) circle (0.1);
    \draw[fill=white] (\j,\sgn*5) circle (0.1);
   }
  }
  \foreach \sx in {-1,1} {
   \foreach \sy in {-1,1} {
    \draw[fill=white] (\sx*4,\sy*4) circle (0.1);
   }
  }
  \fill (0,0) circle (0.13) node[below right] {$0$};
  \node at (-1.6,1.5) {$V$};
  \node[above] at (0,5.6) {cell inside its tile};
  \draw[->] (-1.8,-1.8)--(-3.85,-1.8);
  \node[anchor=west] at (-1.8,-1.8) {$Z$};
  \node at (0,-6.4) {filled boundary sites: $Z$; open sites: $B$};
 \end{scope}
 \begin{scope}[shift={(8.1cm,0.25cm)},x=0.8cm,y=0.8cm]
  \fill[gray!12] (-2.2,-2.2) rectangle (0,-0.03);
  \draw[gray] (-2,-1)--(1,-1);
  \draw[gray] (-1,-2)--(-1,1);
  \draw[gray] (0,-2)--(0,0)--(-2,0);
  \draw[very thick] (0,-1)--(1,-1);
  \draw[very thick] (-1,0)--(-1,1);
  \draw[very thick] (0,-1)--(0,0)--(-1,0);
  \fill (0,-1) circle (0.07);
  \fill (-1,0) circle (0.07);
  \fill (-1,-1) circle (0.07);
  \draw[->] (0,-1.12)--(-0.85,-1.12);
  \draw[->] (-1.12,0)--(-1.12,-0.85);
  \draw[fill=white] (1,-1) circle (0.1) node[right] {$e_1$};
  \draw[fill=white] (-1,1) circle (0.1) node[above] {$e_2$};
  \draw[fill=white] (0,0) circle (0.1)
   node[above right,align=left] {$e_3$\\clipped corner};
  \node[below right] at (0,-1.35) {$(L,L-1)$};
  \node[above left] at (-1,0) {$(L-1,L)$};
  \node[align=center] at (-0.25,-3.0)
   {$Q:\ \operatorname{span}\{e_1+e_3,e_2+e_3\}$\\
    $\ker M:\ \operatorname{span}\{(-1,-1,1)\}$};
 \end{scope}
\end{tikzpicture}
\caption{The clipped cell and a magnified positive corner.  The two
boundary rows couple to a triple of exterior coordinates, leaving exactly
one missing direction.  The inward arrows reach a site outside $Z$;
this geometric fact is used in the transmission argument.}
\label{fig:clipped-cell}
\end{figure}

\subsection{Calibrating the root with a decaying profile}

We choose the root reference value so that the Dirichlet cell has a mode
at energy $E$. For arbitrary real values $v|_J$, write
\[
 A=\operatorname{Adj}_J+\diag_J(v-E).
\]
Whenever $A$ is invertible, define
\begin{equation}
 c(v)=E+p^\top A^{-1}p.
 \label{cell:calibrated-root}
\end{equation}
The Dirichlet matrix on $V$, with root potential $c(v)$, then satisfies
\begin{equation}
 \begin{pmatrix}c(v)-E&p^\top\\p&A\end{pmatrix}
 \begin{pmatrix}1\\-A^{-1}p\end{pmatrix}=0.
 \label{cell:dirichlet-mode}
\end{equation}
Thus invertibility of the punctured matrix $A$ permits resonance of the
full Dirichlet cell. We must choose the interior law so that $A^{-1}$ is
controlled, $c(v)$ stays inside the physical potential interval, and this
root-normalized mode is small at the boundary. The transmission estimates
will later control how it couples to an arbitrary passive exterior.

We first construct a comparison profile $\psi$, normalized by $\psi_0=1$,
and potentials $v^0$ satisfying its local equation. The profile is defined
also on $B$, where it is small but nonzero, so it need not be a Dirichlet
eigenvector. A deformation of this profile will give a polynomial inverse
bound for the punctured matrix at $v^0$. These two properties will control
the exact calibrated mode after the interior values are randomly perturbed.

Choose a rational multiple $\theta$ of $\pi$ and a real phase $\phi$ such
that
\begin{equation}
 |E-4\cos\theta|<h,
 \qquad m_\theta:=\min_{k\in\Z}|\cos(\theta k+\phi)|>0.
 \label{cell:carrier}
\end{equation}
The first choice is possible because the energy is in $(-4-h,4+h)$ and
rational multiples of $\pi$ are dense.  The sequence of cosine values is
periodic, so the second condition excludes only finitely many phases
modulo $2\pi$.  Put
\[
 d=\frac{h-|E-4\cos\theta|}{8}>0.
\]
Define the real sequence $f$ on all integers by
\begin{equation}
 \begin{aligned}
 f_0&=0,\\
 f_k+f_{k+1}
  &=\operatorname{sign}\bigl(\cos(\theta k+\phi)
                  \cos(\theta(k+1)+\phi)\bigr).
 \end{aligned}
 \label{cell:recurrence}
\end{equation}
Recursion in either direction gives $|f_k|\le|k|$.
For small fixed $\nu,c_u>0$, and $|u|\le c_u/L$, set
\begin{equation}
 \begin{aligned}
 \psi_x(u)&=e^{-\nu|x|_\infty}
       \frac{\cos(\theta(x_1+x_2)+\phi)}{\cos\phi}
       e^{u f_{x_1+x_2}}\quad(x\in V\cup B),\\
 v_x^0(u)&=E-\sum_{y\sim x}\frac{\psi_y(u)}{\psi_x(u)}
       \quad(x\in V).
 \end{aligned}
 \label{cell:profile-formula}
\end{equation}
Every neighbor in the last sum belongs to $V\cup B$ and $\psi_0=1$.
Uniformly over the allowed parameters,
\begin{equation}
 c e^{-\nu|x|_\infty}\le |\psi_x(u)|
       \le C e^{-\nu|x|_\infty}.
 \label{cell:profile-size}
\end{equation}
Indeed $|u f_{x_1+x_2}|$ is bounded by a fixed multiple of $c_u$ and the
periodic carrier is bounded away from zero.  For neighbors $x,y$ the
additional ratio factors differ from $1$ by $O(\nu+c_u)$: their radial
exponents differ by at most $\nu$, and
$|u(f_{y_1+y_2}-f_{x_1+x_2})|\le Cc_u$.  At $\nu=u=0$, the cosine
addition formula gives $v_x^0=E-4\cos\theta$.  We therefore fix $\nu,c_u$
small enough, once and for all, that
\begin{equation}
 |v_x^0(u)-(E-4\cos\theta)|\le d,
 \qquad |v_x^0(u)|\le h-7d.
 \label{cell:design-margin}
\end{equation}

\begin{lemma}[Choice of the profile parameter]\label{cell:profile}
There is a fixed exponent $C_2$ such that, for every sufficiently large
$L$, a deterministic $u_L\in[-c_u/L,c_u/L]$ can be chosen for which
the real symmetric Dirichlet matrix
\[
 A^0=\operatorname{Adj}_J+\diag_J(v^0-E)
\]
satisfies $\norm{(A^0)^{-1}}\le L^{C_2}$.  The bounds in
\Cref{cell:profile-size,cell:design-margin} hold for this
choice, with constants independent of $L$.
\end{lemma}

\begin{proof}
During the proof retain $u$ as a variable.  Let
$\Psi(u)=\diag(\psi_x(u):x\in J)$, $W(u)=\Psi(u)^2$, and
$Q(u)=-\Psi(u)A^0(u)\Psi(u)$.  Extend a test vector $\alpha$ by zero
at $0$ and on $B$.  Counting each unoriented edge meeting $J$ once gives
\begin{equation}
 \alpha^*Q(u)\alpha
 =\sum_{\{x,y\}:\{x,y\}\cap J\ne\varnothing}
       \psi_x(u)\psi_y(u)|\alpha_x-\alpha_y|^2.
 \label{cell:weighted-form}
\end{equation}
The sum here contains only edges with the other endpoint in $J\cup\{0\}
\cup B$.  Its weights may have either sign.  Nevertheless,
\Cref{cell:recurrence} gives the exact identity
\[
 \frac{\dd}{\dd u}\bigl(\psi_x(u)\psi_y(u)\bigr)
       =|\psi_x(u)\psi_y(u)|
       \quad(x\sim y).
\]
Thus $Q'(u)$ is a positive weighted Dirichlet form.

For each $x\in J$, choose a nearest-neighbor path to $0$ by decreasing
coordinates in absolute value.  It lies in $V$, has at most $2L$ edges,
and its max radius never exceeds $|x|_\infty$.  Every edge on the path
has weight $|\psi_y\psi_z|\ge c|\psi_x|^2$, by
\Cref{cell:profile-size}.  Weighted Cauchy--Schwarz along the path gives
\[
 |\psi_x|^2|\alpha_x|^2
 \le CL\sum_{\{y,z\}\text{ on the path}}
                  |\psi_y\psi_z||\alpha_y-\alpha_z|^2
 \le CL\,\alpha^*Q'(u)\alpha.
\]
Summing over at most $C L^2$ sites shows
\begin{equation}
 Q'(u)\ge cL^{-3}W(u)\ge L^{-C_1}W(u)
 \label{cell:form-monotonicity}
\end{equation}
for some fixed $C_1$, after increasing the lower size threshold.  One may
take $C_1=4$.  Moreover, all $W(u)$ in the allowed interval are comparable
as quadratic forms, with a fixed comparison constant, and
\[
 -CLW(u)\le W'(u)\le CLW(u).
\]

Choose $C_2>C_1+3$ and put $t=L^{-C_2}$.  For sufficiently large $L$,
\[
 (Q\pm tW)'\ge\tfrac12 L^{-C_1}W.
\]
More explicitly, form comparability and \Cref{cell:form-monotonicity}
give, whenever $u_2-u_1>C'tL^{C_1}$ with a sufficiently large fixed $C'$,
\begin{equation}
 Q(u_2)-tW(u_2)>Q(u_1)+tW(u_1).
 \label{cell:inertia-comparison}
\end{equation}
Let $\lambda_k(u)$ be the $k$th eigenvalue, in increasing order, of
$-A^0(u)$.  Congruence by the invertible real diagonal matrix $\Psi(u)$
preserves the numbers of negative, zero, and positive eigenvalues.
Consequently, if $\lambda_k(u_1)\ge-t$, the $k$th eigenvalue of
$Q(u_1)+tW(u_1)$ is nonnegative.  The strict inequality
in \Cref{cell:inertia-comparison} then makes the $k$th eigenvalue of
$Q(u_2)-tW(u_2)$ positive, so $\lambda_k(u_2)>t$.

It follows that each set
$\{u:|\lambda_k(u)|\le t\}$ has diameter at most $C'tL^{C_1}$.
No monotonicity assertion about the individual $\lambda_k(u)$ is needed.
The union of these sets, over $|J|=O(L^2)$ indices, has length at most
$C L^{C_1+2-C_2}=o(L^{-1})$, whereas the available interval has length
$2c_u/L$.  Choose $u_L$ outside this union.  For definiteness, one may
choose the first rational point in a fixed enumeration that belongs to
the open complement.  Then
$\dist(0,\operatorname{spec}A^0)>L^{-C_2}$, proving the inverse bound.
\end{proof}

From now on $u=u_L$ is fixed and suppressed in the notation.  Set
\begin{equation}
 \epsilon=e^{-\nu L},\qquad \sigma=\epsilon^2.
 \label{cell:epsilon}
\end{equation}
In particular, $|\psi_x|$ is comparable to $\epsilon$ on $B$ and on
$Z^+$, with constants independent of $L$. We have chosen the port strength
$\sigma$ as the square of this designed boundary-profile scale. The exact
calibrated mode will satisfy the polynomially lossy comparison in
\cref{cell:residual}.

\subsection{The branch sampler}

We first describe the boundary conditionings used by the sampler.  For
each corner triple use the fixed basis in \Cref{cell:corner-basis} and put
\begin{equation}
 \mathsf D_j=\diag\left(\frac{(r_j)_i}{n_i}:1\le i\le3\right)
 \quad(j=1,2).
 \label{cell:box-directions}
\end{equation}
These matrices satisfy
\[
 \mathsf D_jn=r_j,\qquad \tr\mathsf D_j=0,
 \qquad n^\top\mathsf D_jn=0.
\]
Thus a common shift changes the real diagonal in the missing direction
$n$, while a shift along $\mathsf D_j$ changes its coupling to $r_j$
without changing that diagonal. The three diagonal directions
$I,\mathsf D_1,\mathsf D_2$ are linearly independent:
multiplying a vanishing linear combination by $n$ gives a vanishing
linear combination of the orthonormal vectors $n,r_1,r_2$.
Thus centered coordinates $(s,t_1,t_2)$ specify physical potentials by
\begin{equation}
 \diag(v)=\diag(\bar v)+sI+t_1\mathsf D_1+t_2\mathsf D_2.
 \label{cell:box-coordinates}
\end{equation}
The map has a fixed nonzero Jacobian.  In fact its three columns, as
vectors of diagonal entries, are orthogonal and have length $\sqrt3$.
Use this same coordinate convention at all levels, including at exactly
coincident triples.

A prescribed corner box is the image under \Cref{cell:box-coordinates}
of a product of three centered intervals.  Fix an exponent $\kappa$.
The boxes allowed below have all three side lengths at least
$c(1+L)^{-\kappa}$ and at most $C$, lie at distance at least
$c(1+L)^{-\kappa}$ from the boundary of $(-h,h)^3$, and keep each physical
coordinate in a single polynomial piece of the fixed artificial profile
$\vartheta$.  Constants may depend on $\kappa$; in the hierarchy this
exponent is fixed by $K_0$.  Uniform measure in these box coordinates
is also uniform measure in the physical parallelepiped.  Because
$\vartheta$ is positive in the interior and comparable to squared
distance from the support endpoints, every nonzero artificial profile
$a_i\in\{1,\vartheta\}$ on a permitted box obeys
\begin{equation}
 a_i(v_i)\ge c(1+L)^{-2\kappa}.
 \label{cell:box-profile-lower}
\end{equation}
Increasing the exponent absorbs any fixed comparison constants.

Choose $D_1>C_2+2$.  A shell-depth parameter $D_2$ will be fixed after
the polynomial thresholds in the boundary and zero-jitter transmission
estimates have been fixed.  Finally choose $D_3$ sufficiently large
depending on $D_1,D_2$ and these thresholds.  On $J$ sample
\begin{equation}
 v_x=v_x^0+r\one_{Z^+}(x)+\zeta_x,
 \qquad r\sim\operatorname{Unif}[-L^{-D_1},L^{-D_1}],
 \label{cell:interior-sampler}
\end{equation}
where the $\zeta_x$ are independent centered uniforms of half-width
\begin{equation}
 w_x=
 \begin{cases}
 L^{-D_3},& L-|x|_\infty\le D_2\log L+5,\\
 \epsilon L^{-D_3},& L-|x|_\infty>D_2\log L+5.
 \end{cases}
 \label{cell:jitter-widths}
\end{equation}
They are independent of $r$ and all boundary sampling.

For these interior draws, use the matrix $A$ and the calibrated value
$c(v)$ from \cref{cell:calibrated-root}. At the root, put an open cell
port $a$ of strength $\sigma$, and use $c(v)$ as its real reference
potential. To define the hidden true potential, fix a smooth cutoff $\chi:\R\to[0,1]$ with $\chi=1$ on
$[-d,d]$ and $\supp\chi\subset[-2d,2d]$, and sample
\begin{equation}
 \begin{gathered}
 v_0=c(v)+\delta,\qquad
 \PP(\delta\in\dd t\mid v|_J)
   =\frac{1}{Z_\sigma}\chi(t)\frac{\sigma}{\pi(t^2+\sigma^2)}\,\dd t,\\
 Z_\sigma=\int_{\R}\chi(t)\frac{\sigma}{\pi(t^2+\sigma^2)}\,\dd t.
 \end{gathered}
 \label{cell:hidden-root}
\end{equation}
This is exactly the cutoff Cauchy law used in \Cref{mat:cutoff}.
The open reference $c(v)$ does not depend on the hidden draw $\delta$.

The boundary law is as follows.  On each corner triple sample uniformly
in its prescribed box, independently between the four triples.  On
every single side site $x$ sample
\begin{equation}
 v_x=\xi_x+s_0,\qquad
 \xi_x\sim\operatorname{Unif}[-h/4,h/4],\qquad
 s_0\sim\operatorname{Unif}[-h/4,h/4],
 \label{cell:side-law}
\end{equation}
using independent $\xi_x$ and one common $s_0$.  These draws are
independent of the triples and the interior draws.  Sites of the tile
outside $V\cup B$ retain their independent uniform $[-h,h]$ laws,
independently of all draws just described.  Apart from the common shifts
$r,s_0$ and the dependence of the root center $c(v)$ on $v|_J$, all
sampling is independent.

\begin{proposition}[Properties of the branch law]\label{cell:branch}
For sufficiently large $L$, the above choices have the following
properties, uniformly over permitted corner boxes:
\begin{enumerate}[label=(\roman*)]
 \item Every interior draw has
 \begin{equation}
  \norm{A^{-1}}\le2L^{C_2},\qquad
  \norm{-A^{-1}p-\psi|_J}\le C L^{C_2+1}\epsilon.
  \label{cell:residual}
 \end{equation}
 The exponents in these bounds are independent of $D_2$, provided $D_3$
 is subsequently chosen sufficiently large.
 \item All true potentials on $V$ lie in a fixed compact subinterval
 of $(-h,h)$.  The root center satisfies
 \begin{equation}
  |c(v)-v_0^0|\le C L^{C_2+1}\epsilon.
  \label{cell:root-calibration}
 \end{equation}
 The plateau of $\vartheta$ can therefore be fixed to contain every
 true value on $V$ and every single side value in $B$, for all the cells
 used.  The artificial absorptions on these sites are constant as
 $r$, the jitters, the common side shift, and the hidden root draw vary.
 \item Relative to the true product law on the tile conditioned on
 these same four corner boxes, the branch has a density
 $\mathcal L$ satisfying
 \begin{equation}
  0\le\mathcal L\le
       \exp\bigl(C L^3+C L^2\log L\bigr),
  \qquad \int\mathcal L\,\dd\PP_{\rm true,boxes}=1.
  \label{cell:likelihood}
 \end{equation}
 The corner-box conditional laws and the laws outside $V\cup B$ are
 unchanged.
\end{enumerate}
Artificial probes may have any prescribed mask of absent, constant, or
profile-dependent strengths, as in \Cref{mat:probes}.  For subsequent cell
estimates their common scale is restricted to
$0\le\tau\le e^{-c_*L}$, with a fixed $c_*>2\nu$.
\end{proposition}

\begin{proof}
The perturbation from $A^0$ is diagonal and has norm at most
$L^{-D_1}+L^{-D_3}$.  Taking $D_3>D_1$, \Cref{cell:profile} and a
Neumann series give the first bound in \Cref{cell:residual}.
The profile equation on $J$ gives the exact residual identity
\[
 A^0\psi|_J+p=-M\psi|_B.
\]
The right-hand side has norm at most $C\sqrt L\epsilon$.  On $Z^+$,
$|\psi_x|\le C\epsilon$, so the additional common-shift residual has
norm at most $CL\epsilon L^{-D_1}$.  On the jitter shell,
\[
 |\psi_x|\le C\epsilon L^{\nu D_2},
\]
where the extra depth $5$ has been absorbed in $C$.  Its jitter residual
is bounded by $CL\epsilon L^{\nu D_2-D_3}$.  Off that shell, the
half-width already has a factor $\epsilon$, and $|\psi_x|\le C$,
giving an upper bound $CL\epsilon L^{-D_3}$.
Thus $D_3>\nu D_2+2$ ensures
\[
 \norm{A\psi|_J+p}\le CL\epsilon.
\]
Multiplication by $A^{-1}$ proves the second bound in
\Cref{cell:residual}.  Notice that the resulting exponent $C_2+1$
was fixed before $D_2$.

Since $\psi_0=1$, the profile equation at $0$ reads
$v_0^0=E-p^\top\psi|_J$.  Subtracting this from the definition of $c(v)$
proves \Cref{cell:root-calibration}.  Increase the minimum cell size so
that all interior perturbations and the last calibration error are at
most $d$.  By \Cref{cell:design-margin}, $|v_x|\le h-6d$ on $J$,
$|c(v)|\le h-6d$, and $|v_0|\le h-4d$ after the hidden draw.
The side values lie in $[-h/2,h/2]$.  Choose the plateau to contain
these compact intervals.  This choice uses only the original design
margin $d$, and is therefore available before the large width and shell
exponents are fixed.  Each present artificial profile on these sites is
then identically $1$, proving the constancy assertion, also at the root.

For the density estimate, first condition on $r$ and $s_0$.  Each jitter
has density at most $(2w_x)^{-1}$ in its physical coordinate, hence
likelihood at most $h/w_x$ relative to a uniform $[-h,h]$ coordinate.
For small $\sigma$, $Z_\sigma\ge1/2$, so the conditional root density
is at most $C/\sigma$.  Its likelihood relative to the true uniform
coordinate is at most $C_h/\sigma$.  Each single side coordinate has
likelihood at most $4$.  Multiplying these bounds before averaging the
common shifts gives
\[
 \log\mathcal L
 \le C|J|+D_3|J|\log L+
       \nu L\,\#\{x\in J:w_x=\epsilon L^{-D_3}\}
       +2\nu L+C|B|
 \le C L^3+C L^2\log L.
\]
An upper bound of this form also holds when some individual displayed
logarithms are negative.  The four corner factors cancel against the
identical box conditionings in the reference law.  The outside factors
cancel as well.  All sampler laws are normalized, so integration gives
the last assertion in \Cref{cell:likelihood}.
\end{proof}

We record explicitly the order of choices used later.  The carrier,
$d,\nu,c_u,C_1,C_2$, the compact cutoff and plateau, and the corner-box
exponent are fixed first.  Next fix $D_1$.  The polynomial threshold in
\Cref{imp:weighted} and the regular inverse thresholds in the zero-jitter
part of \Cref{trans:cell} can then be fixed without $D_2,D_3$.
At depth $L-|x|_\infty=\lfloor D_2\log L\rfloor$,
\Cref{cell:profile-size} gives
\[
 \frac{|\psi_x|}{\epsilon}\ge cL^{\nu D_2}.
\]
Choose $D_2$ large enough that this dominates the already fixed
polynomial errors and reflected-field bounds in that argument.
Finally choose $D_3$ to satisfy both the residual requirements above and
the inverse and path-entry stability requirements of
\Cref{trans:cell}.  Enlarging $D_3$ changes neither the exponents in
\Cref{cell:residual} nor the earlier zero-jitter bounds.  All parameters
are now fixed; only after that do we increase the lower cell-size
threshold.  This order prevents the depth choice from depending on a
polynomial exponent that itself grows with the depth.
